\chapter{价值约束动力学：目标偏置、路径代价与规范更新}
\label{chp:value_gauge_field}

我在本章把原来的“价值规范场”重建为价值约束动力学。智能过程具有方向，并不意味着价值是一种隐藏的物理场；更直接的解释是，系统保存目标、偏好、禁行条件、承诺和风险模型，并据此改变候选状态的排序、更新强度与行动权限。事实模型回答“当前可能是什么”，规范模型回答“在谁的授权、什么时间尺度和哪些边界内，应优先保持或改变什么”。两者彼此耦合，却不能相互替代：事实不会仅凭被厌恶而失真，价值也不会仅凭预测准确而自动正当。本章从状态、目标与可行域开始，分别构造方向偏置、复合价值、非对称路径代价和规范更新，说明数学结果、经验偏好与工程设计各自能够支持到哪里。历史上的风场、势能和规范场语言只保留为直觉提示，不再承担机制证明。

\section{从事实状态到规范状态：方向来自目标与约束}
\label{sec:section_6725}

我先定义需要被区分的对象。设智能系统在时刻 $t$ 的任务状态为 $x_t\in\mathcal X_t$，其中可以同时含有连续属性、离散结构身份和外部接口状态；可执行动作集合为 $\mathcal U_t(x_t)$。事实模型 $P_t(y\mid x,u)$ 描述采取动作后可能观察到什么，规范状态则记为
\begin{equation}
    n_t=(g_t,w_t,C_t,R_t,\kappa_t,\pi_t).
\end{equation}
这里 $g_t$ 是目标或期望结果，$w_t$ 是各评价维度的权重或优先序，$C_t$ 是必须满足的硬约束，$R_t$ 是风险度量，$\kappa_t$ 记录规范来源与授权，$\pi_t$ 是适用范围和过期条件。上述量都不是焦耳能量；它们通常是无量纲评分、带业务单位的损失或逻辑谓词。若不同分量单位不同，就必须先按明确的尺度、阈值或偏序组织，不能直接相加。

给定预测模型和规范状态，一步候选的评价可写为
\begin{equation}
    J_t(x,u;n_t)
    =\mathbb E_{y\sim P_t(\cdot\mid x,u)}[\ell(y,g_t)]
     +R_t(u,y)+\Omega_t(x,u),
\end{equation}
并要求 $(x,u,y)\in C_t$。第一项是任务失真，第二项描述尾部风险或不可逆后果，第三项是资源、延迟或切换代价。这个公式是一种工程模型假设，不是“应然”从概率中自动产生。若 $C_t$ 为空，应报告规范冲突；若候选动作均不可接受，应停机、请求授权或扩展观测，而不是通过降低惩罚强行产生动作。目标下降只是内部指标改善，外部成功仍需执行反馈验证。

在连续且结构固定的区间，可以用目标梯度构造局部方向偏置。设 $x\in\mathbb R^d$，$M_t(x)$ 是对称正定的更新度量，则
\begin{equation}
    b_t(x)=-M_t(x)^{-1}\nabla_x J_t(x;n_t)
\end{equation}
给出当前模型中的候选下降方向。它类似“风向”，但来源是显式目标和度量；目标、输入、约束或度量变化时，$b_t$ 也随之变化。若结构图发生节点合并、承诺撤销或权限改变，就应执行离散编辑和重新评估，不能把不同维度的状态硬接成一条光滑轨迹。连续方向只在局部可行域内有意义，越界后的梯度不提供合法行动。

方向偏置还必须与因果模型分开。因果边 $A\to B$ 表示干预条件下的依赖假设，不表示系统偏好 $B$；价值偏好 $B$ 也不能证明 $A$ 导致 $B$。例如系统希望患者康复，不能因此把某个治疗写成有效原因；它必须使用证据模型预测疗效和副作用。反过来，一个可靠的因果通路也可能因伦理、权限或资源约束而不可执行。原章把价值、语境与因果统一进同一个非阿贝尔势，会掩盖这种逻辑差异。本章保留它们的接口关系，却为每类判断保存独立字段和验证方法。

这种分离还给出了可以被证伪的诊断路径。若预测错误，就应更新事实模型或观测质量；若预测正确而行动仍不合要求，应检查目标解释、约束装载、授权范围或风险阈值；若内部选择合理而外部结果失败，则应检查执行接口和环境扰动。三类失败不能统一归因于“价值强度不足”。例如配送系统已经准确预测道路拥堵，却因把医院急件与普通订单采用同一优先序而延误，问题位于规范状态；若优先序正确而车辆没有收到指令，问题则位于执行链。只有保存这一分解，系统才能把纠错送到真正负责的模块，并用后续数据检验修订是否有效。

HSF-HD 在这里给出一般层原则：智能状态更新由事实可行性与规范可接受性共同筛选。AgentOS 可以把 $n_t$ 的快速部分投影到 $h(t)$，把规范变更事件保存在 $E$，把长期目标和承诺沉淀在 $\Theta$，由 $\Psi$、SPC、TCE 或 Boundary 参与读出、调控和授权；这些是工程实例，不是价值动力学的唯一定义。只要另一种系统能够明确记录目标来源、约束范围、冲突处理和反馈证据，也可以实现同一一般原则。

\section{价值分解：目标、风险、语境与资源耦合}
\label{sec:section_6772}

原章把价值场分解为标量势和矢量势，我改用可审计的计算分解。第一类是终端评价，即轨迹结束或检查点上的结果损失 $\ell_T(x_T,g_T)$；第二类是过程评价，如每一步的风险、延迟、通信和认知负荷；第三类是方向相关的转移代价，如提交比撤销更容易、写入比删除更便宜；第四类是语境门控，它决定哪些评价维度在当前角色、时间窗和授权域内有效。这四类对象的数学性质不同，不能仅因都影响行为就合并成一种“认知力”。

对有限时域 $k=t,\ldots,T$，可将候选策略 $\mu$ 的评价写为
\begin{equation}
    \mathcal J_t(\mu)
    =\mathbb E^\mu\!\left[
      \ell_T(x_T,g_T)+
      \sum_{k=t}^{T-1}
      \bigl(\ell_k(x_k,u_k,g_k)+c_k(x_k,u_k)+\lambda_k r_k(x_k,u_k)\bigr)
    \right],
\end{equation}
同时要求轨迹满足 $C_k$。$c_k$ 可以以秒、调用次数或货币计量，$r_k$ 可以是失效概率、条件风险值或等级评分；若单位不同，权重 $\lambda_k$ 必须承担单位换算或明确只用于排序。目标 $g_k$、权重、约束和输入随时间变化时都已进入求和，不应把它们冻结后仍声称描述在线系统。对无限时域还要声明折扣、平均代价或稳定条件，否则总和可能不存在。

方向性主要来自状态转移，而不是抽象空间的“时间箭头”。设离散状态图的有向边为 $e=(i,j)$，基础计算代价为 $d_e\geq0$，规范偏置为 $a_e$，则有效边代价可定义为 $c_e=d_e+a_e$。允许 $a_{ij}\neq a_{ji}$，就能表达写入/回滚、承诺/撤销、压缩/重建的非对称性。但为避免负环使最短路问题失去下界，需要要求所有可达有向环的总代价非负，或更强地令每条 $c_e\geq\varepsilon>0$。某条逆向边不存在表示接口不允许逆操作，不必夸张为“无限物理阻抗”。恢复困难应由备份、信息丢失和权限模型说明。

语境门控记为 $m_t\in[0,1]^p$，价值向量为 $v_t(x,u)\in\mathbb R^p$。逐分量乘积 $m_t\odot v_t$ 表示当前哪些维度被激活，但门控不应自动取消硬约束。忍痛执行任务可以降低疼痛信号在软目标中的权重，却不能把组织损伤或安全阈值删除；聚焦可以提高任务相关候选的资源份额，却不能令其他责任趋于无穷小。若注意力预算为 $B_t$，资源分配 $r_{t,j}\geq0$ 应满足 $\sum_jr_{t,j}\leq B_t$，这样“耦合强度”才有可检查的上界，而不是任意趋于无穷的神秘荷量。

风险也不能与厌恶强度混为一谈。风险来自结果分布与损失尾部，厌恶参数描述主体对同一风险的态度。两个系统可以面对相同失效概率，却因责任、可恢复性或利益相关者不同采取不同动作；同一系统也可能在紧急状态改变容忍度。应分别记录证据模型 $P_t$、损失 $L_t$ 和风险泛函 $\rho_t(L_t)$，再由规范来源决定阈值。把恐惧直接映射成高电势会丢失“害怕但判断安全”与“平静但风险极高”的区别。

价值分解并不意味着各项彼此独立。一次外部调用可能同时缩短时延、增加隐私暴露并消耗预算，所以模型还需记录共享原因和耦合约束。不过，耦合应通过联合结果变量、条件概率或显式约束表达，而不能用一个总分抹平来源。测试时可以固定预测分布只改变权限，固定权限只改变预算，或固定资源只改变风险阈值，观察选择是否按声明方向变化。这样的反事实测试能发现权重符号写反、门控遗漏和单位缩放错误，也为不同部署环境提供可重复的校准依据。

在工程上，CCM 可以估计候选数量和运行负荷，TCE 调整软目标的增益与探索范围，Boundary 约束外部交换，Offloading 以未来代价换取当前负荷下降。它们共同改变 $\mathcal J_t$ 的不同分量，却不应写成一个不可分辨的总场。逐项记录让系统能够回答：某个动作是因目标收益、资源不足、风险阈值还是权限限制被选择或拒绝。这种可解释分解比统一物理隐喻更适合测试、审计和修订。

\section{复合规范结构：冲突、优先序与可修订承诺}
\label{sec:semantic_structure_of_gauge}

复杂智能体不会只有一个标量价值。设评价向量为 $V(x,u)=(V_1,\ldots,V_p)$，其中可包含安全、真实性、帮助性、个体偏好、组织承诺和资源成本。直接写 $\sum_jw_jV_j$ 隐含可补偿性：一个维度的巨大收益可以抵消另一个维度的严重违反。若安全或合法性不可补偿，就应先用硬约束筛除，再在可行集合上优化软目标。若多个软目标也不可稳定换算，可使用 Pareto 偏序或词典序，而不是虚构一个天然统一的价值单位。

令候选集合为 $\mathcal K_t$，硬约束谓词为 $C_t^q(k)\in\{0,1\}$，则合法集合
\begin{equation}
    \mathcal K_t^{\mathrm{safe}}
    =\{k\in\mathcal K_t:C_t^q(k)=1,\ q=1,\ldots,Q\}.
\end{equation}
如果集合非空，系统再依据偏序 $\preceq_{n_t}$ 选择极小元；若极小元不唯一，可以保留候选、请求更多证据或按可逆性试探。这个过程区分“违反约束”和“偏好较低”。例如在医疗建议中，禁忌症是硬约束，预计等待时间可能是软成本；不能因为时间收益足够大就抵消禁忌。对创作任务，审美维度通常可协商，但来源真实性和用户授权仍可能是不可补偿条件。

价值冲突不是非对易李代数的必然证据，而是规范集合与资源条件不相容。可以构造最小冲突集：若 $C_t$ 整体不可满足，寻找一个不可满足而任意真子集可满足的约束子集，向上游报告具体冲突。若冲突来自不同授权者，还要比较授权范围、时间戳和撤销规则，而不是让数值最大的“价值力”获胜。社会系统中的规范更新尤其需要程序合法性：新偏好不能静默覆盖既有承诺，撤销也应保留事件记录和受影响对象。

承诺可表示为带条件的记录
\begin{equation}
    c=(\mathrm{id},\mathrm{issuer},\mathrm{scope},\mathrm{predicate},
       \mathrm{priority},\mathrm{validity},\mathrm{revocation}).
\end{equation}
更新函数 $U_n(n_t,e_{t+1})$ 接受新证据或授权事件，产生 $n_{t+1}$ 及差异日志。若只是坐标或表述改变，语义等价应由迁移规则保持；若优先序、适用域或授权者改变，则属于真正的规范事件。这样可以区分“换一种视角解释事实”与“改变应该做什么”。前者可能改变读出方式，后者必须触发重新规划、影响分析和必要的外部确认。

情绪和习惯可以进入价值动力学，但它们是状态变量而非证明。情绪 $e_t$ 可以改变注意门控、风险容忍或候选生成，从而形成短时偏置；习惯可以通过重复更新降低某条策略的检索或执行成本。两者都应受边界和反馈校正：愤怒可能放大威胁评分，却不能自动提高威胁事实的可信度；熟练习惯降低执行延迟，却不能证明在新语境仍合适。经验主张应由行为数据验证，不能仅由“场曲率”或“相位”措辞推出。

规范更新还要处理时间一致性。短期偏好可以改变，但已经据此影响他人的承诺不能被无痕重写；未来自我也不必被过去配置永久锁定。我的做法是把偏好版本、依赖该版本作出的决策以及可撤销条件同时保存。更新发生后，系统先计算受影响集合，再区分可自动迁移、需要重新确认和必须维持的事项。这样，稳定性来自有条件的连续承诺，而不是假定价值永恒不变；适应性来自受控修订，也不是让最近一次情绪覆盖全部历史责任。

在 HSF-HD 中，复合规范结构与世界结构、分布表示共同构成状态，但三者职责不同：世界结构保存对象与关系，分布表示保存不确定性，规范结构保存目标、约束和授权。AgentOS 的 $\Psi$ 可以提供长期自我相关优先序，SPC 形成当前自我状态读出，TCE 调节软增益，$E$ 保存承诺事件，$\Theta$ 保存持久规则；任何工程实现都必须防止自我偏好绕过事实校验和外部边界。一般理论只要求这些职责可被区分、组合和审计。

\section{方向相关路径代价：连续模型、离散实现与适用边界}
\label{sec:section_5960}

原章用 Randers 度量表达“顺风与逆风”代价不同，这个数学工具可以保留，但必须补足正定条件并重新解释。设 $x$ 位于光滑状态区域，$v\in T_x\mathcal X$ 是无穷小状态变化，$G(x)$ 是对称正定矩阵，$b(x)$ 是一形式。定义
\begin{equation}
    F(x,v)=\sqrt{v^{\mathsf T}G(x)v}+b(x)^{\mathsf T}v.
\end{equation}
若对偶范数满足 $\lVert b(x)\rVert_{G^{-1}}<1$，则由 Cauchy--Schwarz 不等式，$b^{\mathsf T}v\geq-\lVert b\rVert_{G^{-1}}\sqrt{v^{\mathsf T}Gv}$，所以任意 $v\neq0$ 都有 $F(x,v)>0$。这排除了原章所谓代价任意变成负数的情况。$F(x,v)$ 一般不等于 $F(x,-v)$，因此它能表示方向相关的局部计算代价，但不是焦耳能耗，也不是价值真实性。

对绝对连续路径 $\gamma:[0,1]\to\mathcal X$，定义累计代价
\begin{equation}
    L(\gamma)=\int_0^1F(\gamma(s),\dot\gamma(s))\,ds.
\end{equation}
第一项刻画在表示度量下的基础变化规模，第二项刻画目标、习惯或承诺造成的方向偏置。若 $b=\nabla\varphi$ 是精确一形式，第二项只由端点差 $\varphi(\gamma(1))-\varphi(\gamma(0))$ 决定；若不是，路径历史会影响累计代价。即便如此，最小 $L$ 路径也只是在给定模型中的最省代价路径，不保证逻辑证明有效、伦理可接受或外部执行成功。硬约束必须通过可行路径集合单独施加。

连续模型不适合所有状态变化。工具调用、身份合并、权限授予和承诺撤销是离散事件，宜用有向图 $H_t=(S_t,E_t)$ 表示。每条边保存基础代价、方向偏置、前置条件、可逆性和失败恢复方案。路径代价是边代价之和，图搜索只在满足前置条件的子图上运行。若边权随目标和环境在线变化，算法必须使用当前版本，并在执行后重新规划；不能把初始最短路当作整个过程的固定保证。若存在负边，应使用能处理负权的算法并排除可达负环；更稳妥的规范系统通常直接保持非负边代价。

以“忍住即时奖励完成长期任务”为例，基础状态距离并没有改变；改变的是短期奖励、长期目标、失败风险和注意预算的权衡。系统可能通过提高长期承诺的优先级、降低诱因候选的曝光或把诱因移出环境来改变路径。以逆向工程为例，正向生成与从结果恢复过程不同，是因为信息丢失、搜索分支和接口权限不同，而不是抽象空间施加了无限风阻。具体机制解释比方向隐喻更能预测何时可以借助日志、检查点或外部记忆降低逆向代价。

方向偏置还需要安全边界。若软目标足以把系统持续推向可行域边缘，应使用投影、屏障函数或运行时监控阻止越界；若目标变化过快造成振荡，应限制增益、引入滞回或按时间尺度分层更新。连续控制需声明步长与稳定条件，离散执行需声明事务提交和回滚条件。AgentOS 中，TCE 可以调节增益，CCM 监测负荷和冲突，Boundary 限制交换与执行，Offloading 改变未来路径成本；这些机制使“方向性”成为可观察、可干预的运行变量。

离散化时，若用步长 $\Delta t$ 近似连续更新，就应验证局部误差和约束裕量是否仍在允许范围内；减小步长通常提高近似精度，却会增加计算与通信次数。结构事件则不能靠无限减小步长解决，而要在事务边界上重新构造状态空间。因而连续路径模型负责同一表示结构内的渐进调整，有向图和事件日志负责结构改变，两者通过版本号、前置条件和读出接口衔接。

因此，本章不再声称目的论被写进宇宙几何，而主张智能系统必须把事实预测、规范状态和方向相关代价联合起来。HSF-HD 作为智能的一般状态动力学，描述目标与约束怎样改变状态转移的选择分布、路径成本和稳定性；AgentOS 再把这些关系落实为记忆、调度、边界与反馈接口。这个层级既保留了“同一路程因方向不同而代价不同”的洞见，也避免用物理名称遮蔽价值来源、授权责任和验证边界。
